\section{Absorber and thermodynamic model}\label{sec:model-framework}
\subsection{Physical system and model structure}
% M1: physical domain
The model describes steady countercurrent absorption in one packed bed.
Vapor flows upward from $z=0$ to $z=H$, while lean solvent enters at $z=H$ and flows downward; $H$ is packed height and $A$ is column cross-sectional area.
The apparent liquid components are \COtwo, MEA and water, and the wet gas contains \COtwo, water, nitrogen and oxygen.
Only \COtwo{} and water transfer between phases.
Two models of this system are evaluated.
The nine-species model integrates the column balances at constant pressure and carries the sensitivity, operating and computational-cost studies.
The twelve-state model (\cref{subsec:twelve-state}) solves conserved cell balances with a hydraulic pressure drop and leads the capture and temperature comparison.

% M2: calculation sequence
At each axial position, reactive equilibrium and transport calculations supply the fluxes used in the column balances (\cref{fig:absorber-calculation-sequence,tab:thermodynamic-responsibility}).
ePC-SAFT supplies every bulk-phase thermodynamic property in these balances: phase and reaction equilibrium, liquid and vapor densities, phase enthalpies and heat capacities, and the \COtwo{} and water fugacities.
With the MEA parameter set of \cref{sec:selected-reactive-inputs}, liquid temperature, pressure and apparent composition determine bulk speciation, density, enthalpy and the liquid \COtwo{} and water fugacities.
A wet four-species PC-SAFT vapor with published parameters supplies the water fugacity coefficient, vapor density and vapor enthalpy, and a reduced neutral vapor composition supplies the vapor \COtwo{} fugacity coefficient.
Surface tension, viscosity, diffusivity, reaction kinetics and packing performance remain correlations, and the \NtwoO{}-analogy Henry correlation of \cref{sec:properties} only generates the starting profile.
These quantities enter conventional enhancement and transfer resistances, and the resulting material and enthalpy sources determine the phase-flow and temperature derivatives.

\begin{figure}[pos=htbp]
\centering
\begin{tikzpicture}[node distance=6mm and 9mm]
  \node[cse stage] (state) {Phase state};
  \node[cse detail,right=of state] (state-detail)
    {Flows, apparent composition, temperatures and pressure};
  \node[cse stage,below=12mm of state] (thermo) {Reactive liquid state};
  \node[cse detail,right=of thermo] (thermo-detail)
    {Five reactions and EOS activities give nine species, density, enthalpy and liquid CO$_2$ and water fugacities};
  \node[cse stage,below=12mm of thermo] (vapor) {Vapor state};
  \node[cse detail,right=of vapor] (vapor-detail)
    {Reduced CO$_2$/trace-MEA/water EOS $\phi^v$ with wet-gas $y_{\mathrm{CO_2}}$; wet-gas EOS water fugacity, density and enthalpy};
  \node[cse stage,below=12mm of vapor] (properties) {Transport properties};
  \node[cse detail,right=of properties] (properties-detail)
    {Surface tension, viscosity, diffusivities and heat-transfer coefficient from correlations};
  \node[cse stage,below=12mm of properties] (transport) {Interphase sources};
  \node[cse detail,right=of transport] (transport-detail)
    {Conventional enhancement, fixed-bulk-activity film resistance, mass and heat transfer};
  \node[cse stage,below=12mm of transport] (balances) {Column balances};
  \node[cse detail,right=of balances] (balances-detail)
    {Component-flow and temperature derivatives with mixed inlet boundaries};
  \node[cse solver,below=12mm of balances] (solver) {Collocation update};
  \node[cse detail,right=of solver] (solver-detail)
    {Update axial states and mesh using differential and boundary residuals};
  \foreach \from/\to in {state/thermo,thermo/vapor,vapor/properties,properties/transport,transport/balances,balances/solver}
    \draw[cse arrow] (\from) -- (\to);
  \foreach \name in {state,thermo,vapor,properties,transport,balances,solver}
    \draw[cse link] (\name.east) -- (\name-detail.west);
  \draw[cse arrow] (solver.west) -- ++(-9mm,0) |- (state.west);
\end{tikzpicture}

\caption{Coupled nine-species absorber calculation, from reactive speciation and phase driving forces through conventional enhancement and interphase sources to the column balances; feedback updates the axial state. A reduced CO$_2$/trace-MEA/water composition supplies the vapor CO$_2$ fugacity coefficient, and the wet four-species vapor supplies the water fugacity, density and enthalpy.}
\label{fig:absorber-calculation-sequence}
\end{figure}

\begin{table}[pos=htbp]
\centering\small
\caption{Quantities supplied by each part of the nine-species absorber model.}
\label{tab:thermodynamic-responsibility}
\begin{tabularx}{\textwidth}{@{}>{\raggedright\arraybackslash}p{0.24\textwidth}>{\raggedright\arraybackslash}X>{\raggedright\arraybackslash}X@{}}
\toprule
Calculation & Supplied quantities & Assumption or boundary \\
\midrule
Reactive liquid equilibrium and ePC-SAFT & Nine-species amounts, activities, density, enthalpy, apparent partial enthalpies and heat capacity; liquid CO$_2$ and water fugacities & Five reaction constants, conserved components and electroneutrality \\
Reduced neutral vapor EOS & CO$_2$ fugacity coefficient at total pressure & N$_2$/O$_2$ excluded from EOS composition; original wet-gas fraction used in $f^v=y_{\COtwo}\phi^vP$ \\
Wet four-species vapor PC-SAFT & Water fugacity coefficient, density, enthalpy, partial enthalpies and heat capacity; inlet water saturation & Published pure-component parameters, zero binary corrections and ideal-gas heat capacities \\
Enhancement-factor film & Enhancement and combined transfer resistance & $H_{\mathrm{bulk}}=f^\ell_{\COtwo}/C^{\ell,\mathrm{free}}_{\COtwo}$; bulk activity coefficient fixed across film \\
Correlations & Surface tension, viscosity, diffusivities, heat-transfer coefficient, kinetics and packing performance & Source ranges in \cref{tab:transport-applicability} \\
\bottomrule
\end{tabularx}
\end{table}

\subsection{Reactive equilibrium and phase driving forces}\label{subsec:thermodynamics}
% M3: species and conservation
Five reactions distribute the conserved apparent MEA--\COtwo--water composition among nine liquid species.
The liquid EOS and chemical-equilibrium equations are solved together at the liquid temperature and pressure of each axial position.
The nine true species, in the order used below, are \(\COtwo\), \(\MEA\), \(\HtwoO\), \(\MEAH\), \(\MEACOO\), \(\HCOthree\), \(\ce{CO3^{2-}}\), \(\ce{H3O+}\) and \(\ce{OH-}\).
The independent reactions are
\begin{align}
2\HtwoO &\rightleftharpoons \ce{H3O+}+\ce{OH-}, \tag{R1}\label{eq:reactive-r1}\\
\COtwo+2\HtwoO &\rightleftharpoons \HCOthree+\ce{H3O+}, \tag{R2}\label{eq:reactive-r2}\\
\HCOthree+\HtwoO &\rightleftharpoons \ce{CO3^{2-}}+\ce{H3O+}, \tag{R3}\label{eq:reactive-r3}\\
\MEACOO+\HtwoO &\rightleftharpoons \MEA+\HCOthree, \tag{R4}\label{eq:reactive-r4}\\
\MEAH+\HtwoO &\rightleftharpoons \MEA+\ce{H3O+}. \tag{R5}\label{eq:reactive-r5}
\end{align}
With products assigned positive stoichiometric coefficients, the equilibrium and conservation conditions are
\begin{align}
\sum_i\nu_{ir}\ln a_i^{\mathrm{EOS}}&=\ln K_r^{\mathrm{EOS}}(T,P),&
\mathbf B\mathbf n&=\mathbf B\mathbf n_0,\\
\mathbf M^{\mathsf T}\mathbf n&=\mathbf M^{\mathsf T}\mathbf n_0,&
\mathbf z^{\mathsf T}\mathbf n&=0.
\end{align}
Here \(\mathbf n\) contains true-species amounts, \(\mathbf n_0\) contains the apparent \(\COtwo\), MEA and water feed amounts with zero initial ionic amounts, \(\mathbf M\) contains formula-based molar masses, and \(\mathbf z=(0,0,0,1,-1,-1,-2,1,-1)^{\mathsf T}\).
The carbon and nitrogen balances use
\begin{equation}
\mathbf B=
\begin{bmatrix}
1&2&0&2&3&1&1&0&0\\
0&1&0&1&1&0&0&0&0
\end{bmatrix}.
\end{equation}
Mass conservation supplies the remaining independent composition constraint.

% M4: activities and standard states
Reaction equilibrium is evaluated on the molality activity basis with temperature-dependent constants.
The source standard state uses pure water and infinite-dilution solute molality, with \(m^\circ=\qty{1}{mol\,kg^{-1}}\).
Here $m_i^{\mathrm{mol}}$ is solute molality in moles per kilogram of solvent and $\gamma_i^{(m)}$ is its dimensionless molality-basis activity coefficient.
Source-basis solute activities are \(a_i^{\mathrm{src}}=\gamma_i^{(m)}m_i^{\mathrm{mol}}/m^\circ\); water uses its pure-solvent reference.
Reaction constants and chemical potentials are converted consistently to the EOS reference convention used in the equilibrium equations above.
Hereafter $K_r$ denotes the effective $K_r^{\mathrm{EOS}}(T,P)$, including its pressure-dependent reference conversion; temperature dependence is discussed at the fixed pressure of each column calculation.
\Cref{sec:selected-reactive-inputs} distinguishes the temperature correlations from these EOS-basis constants.

All five constants are fixed literature correlations on their source bases: R1--R3 from Austgen et al.\ \cite{Austgen1991}, R4 from Tong et al.\ \cite{Tong2012} and R5 from Bates and Pinching \cite{batesAcidicDissociationConstant1951}.
None was fitted with the parameter set.
The coefficients, source conventions and calculation domains are given in \cref{sec:selected-reactive-inputs}; their common temperature range is \qtyrange{293.15}{393.15}{K}.

The coupled liquid EOS supplies normalized true-species mole fractions and molar density, giving
\begin{equation}
C_i^{\mathrm{true}}=\rho_{\mathrm{EOS}}^{\liq}x_i^{\mathrm{true}}.
\end{equation}
These free and ionic concentrations enter reaction enhancement.
The same EOS density supplies the hydraulic calculation and the apparent-MEA concentrations used for diffusivity.


% M5: EOS physical contributions
The residual Helmholtz-energy model represents the physical interactions entering activities and fugacity coefficients.
For the liquid phase, the MEA electrolyte state activates the ionic terms.
In the vapor the same expressions apply without the electrolyte terms, because the vapor contains only neutral species.
For each phase state $(T,P,\mathbf{x})$, the EOS solves for the phase density and evaluates the residual Helmholtz energy as $\tilde{a}^{res}=\tilde{a}^{hc}+\tilde{a}^{disp}+\tilde{a}^{assoc}+\tilde{a}^{DH}+\tilde{a}^{Born}$.
The hard-chain term accounts for excluded-volume and chain-connectivity effects, the dispersion term accounts for attractive interactions, the association term accounts for site-specific association such as hydrogen bonding, the Debye--H\"uckel term accounts for long-range ionic screening, and the Born term accounts for ion-solvation work.
The five contributions are
\begin{align}
\tilde{a}^{hc}
&= \bar{m}\tilde{a}^{hs}
-\sum_i x_i(m_i-1)\ln g_{ii}^{hs}(\sigma_{ii})
&& \text{dimensionless}\label{eq:epcsaft-hc}\\
\tilde{a}^{disp}
&= -2\pi\widehat\rho I_1(\eta,\bar{m})\overline{m^2\epsilon\sigma^3}
-\pi\widehat\rho\bar{m}C_1I_2(\eta,\bar{m})\overline{m^2\epsilon^2\sigma^3}
&& \text{dimensionless}\label{eq:epcsaft-disp}\\
\tilde{a}^{assoc}
&= \sum_i x_i\sum_{A_i}
\left(\ln X^{A_i}-\frac{X^{A_i}}{2}+\frac{1}{2}\right)
&& \text{dimensionless}\label{eq:epcsaft-assoc}\\
\tilde{a}^{DH}
&=-\frac{\kappa e^{2}}{12\pi\varepsilon_{0}\varepsilon_{r,\mathrm{bulk}}k_{B}T}
\sum_i x_i z_i^2 \chi_i
&& \text{dimensionless}\label{eq:epcsaft-dh}\\
\tilde{a}^{Born}
&=
-\frac{e^2}{4\pi\varepsilon_0 k_B T}
\sum_i x_i z_i^2
\sum_{\delta\in\{\mathrm{bulk},\mathrm{ion}\}}
\left(1-\frac{1}{\varepsilon_{r,\delta}}\right)
D_i^{\delta}
&& \text{dimensionless}\label{eq:epcsaft-born}.
\end{align}
Here $m_i$ is the segment count of species $i$, $\bar{m}=\sum_i x_i m_i$ is the mean segment number and $\widehat\rho=10^{-30}N_A\rho_{\mathrm{mol}}$ is molecular number density in \unit{\angstrom^{-3}} when $\rho_{\mathrm{mol}}$ is in \unit{mol\,m^{-3}}.
The hard-sphere residual Helmholtz energy is $\tilde a^{hs}$, and $g_{ii}^{hs}$ is the hard-sphere contact radial distribution function.
The functions $I_1$ and $I_2$ are the first- and second-order dispersion integrals; $C_1$ is the hard-chain compressibility correction in the cited PC-SAFT formulation.
Pair diameters and energies obey $\sigma_{ij}=(\sigma_i+\sigma_j)/2$ and $\epsilon_{ij}=(1-k_{ij})\sqrt{\epsilon_i\epsilon_j}$.
The dispersion moments use segment diameters numerically in \unit{\angstrom} and reduced pair energies:
\begin{equation}
\overline{m^2\epsilon^n\sigma^3}
=\sum_i\sum_j x_i x_j m_i m_j
\left(\frac{\epsilon_{ij}/k_B}{T}\right)^n\sigma_{ij}^3,
\qquad n=1,2.
\end{equation}
Thus the density--moment products are dimensionless.
Further, $\eta$ is the packing fraction, $X^{A_i}$ is the fraction of nonbonded association sites, $\kappa$ is the inverse Debye screening length, and $\chi_i$ is the finite-ion-size screening factor.
The Born sum covers the bulk-solvent and ionic-shell regions, with relative permittivities $\varepsilon_{r,\delta}$ and reciprocal-diameter terms $D_i^{\delta}$ defined in \cref{sec:selected-reactive-inputs}.
The Born contribution follows Figiel et al.\ \cite{Figiel2025}, including its solvation-shell and dielectric-saturation (SSM+DS) terms.
The remaining auxiliary definitions, derivative expressions and mixing rules follow the cited PC-SAFT and ePC-SAFT formulations \cite{Gross2001,Cameretti2005,Bulow2019,Bulow2020,Bulow2021a}; the model choices and numerical inputs are given in \cref{sec:selected-reactive-inputs}.

% M6: actual phase driving forces
Interphase \COtwo{} transfer uses liquid fugacity from the reactive state and a vapor fugacity coefficient evaluated on a reduced neutral composition.
The vapor EOS excludes nitrogen and oxygen and normalizes CO$_2$, a numerical MEA trace and water:
\begin{align}
S_v&=\max(y_{\COtwo},\delta_x)+\delta_x+\max(y_{\HtwoO},\delta_x),\\
\widetilde{\mathbf y}&=\frac{1}{S_v}
\left(\max(y_{\COtwo},\delta_x),\delta_x,\max(y_{\HtwoO},\delta_x)\right),
\qquad\delta_x=10^{-12},\label{eq:reduced-vapor}\end{align}
where $\mathbf y$ is the full wet-gas composition and $\delta_x$ is a numerical MEA trace.
The coefficient evaluated at $T^\vap$, total pressure $P$ and $\widetilde{\mathbf y}$ is multiplied by the original wet-gas CO$_2$ fraction:
\begin{align}
f_{\COtwo}^{\vap}&=y_{\COtwo}\phi_{\COtwo}^{\vap}(T^\vap,P,\widetilde{\mathbf y})P,\\
f_{\COtwo}^{\liq}&=x_{\COtwo}^{\mathrm{true}}\phi_{\COtwo}^{\liq}(T^\liq,P,\mathbf x^{\mathrm{true}})P.
\label{eq:co2-fugacity}
\end{align}
Both fugacities are in pascals; the liquid fraction is the true-species mole fraction of unreacted molecular \COtwo{}, not apparent absorbed carbon.

Water fugacity is evaluated on the full phase compositions:
\begin{align}
f^{\:\vap}_{ \HtwoO} &= y_{\HtwoO}\phi^{\vap}_{\HtwoO}(T^\vap,P,\mathbf y)P && \left(\unit{Pa} \right)\\
f^{\:\liq}_{ \HtwoO} &= x^{\mathrm{true}}_{\HtwoO}\phi^{\liq}_{\HtwoO}(T^\liq,P,\mathbf x^{\mathrm{true}})P && \left(\unit{Pa} \right)
\end{align}
The liquid coefficient comes from the nine-species reactive state.
The vapor coefficient comes from a wet four-species PC-SAFT vapor of \COtwo{}, water, nitrogen and oxygen with published pure-component parameters \cite{Gross2001,Schick2023} and zero binary corrections; the same vapor state supplies vapor density, enthalpy and heat capacity.
The vapor--liquid fugacity difference supplies the driving force, while transfer coefficients, interfacial area and reaction enhancement determine the flux.

% M7: parameter origin
The liquid calculation uses the MEA parameter set of the authors' amine-thermodynamics study (\cref{sec:selected-reactive-inputs}).
Water, the two MEA-derived ions and their binary interactions were fitted jointly to vapor--liquid equilibrium, \COtwo{} solubility, speciation, absorption-heat, density and heat-capacity data; the remaining inputs are literature values or separate fits listed in \cref{sec:selected-reactive-inputs}.
No thermodynamic parameter was fitted to the NCCC cases, so the column comparisons are predictions of the parameter set.
Because liquid densities are in the fit objective, the density comparison of \cref{sec:properties} uses only rows outside it.
Ideal-gas heat capacities and formation enthalpies of \COtwo{} and water \cite{chaseNISTJANAFThermochemicalTables1998,internationalassociationforthepropertiesofwaterandsteamRevisedReleaseIAPWS2018} and of MEA \cite{Zhang2011} complete the caloric description of both phases.
Reaction enthalpies follow from the temperature derivatives of the five equilibrium constants.

\subsection{Hydraulics, transfer coefficients and enhancement}\label{subsec:transport}
% M8: hydraulic dependencies
Flows and mixture properties at each axial position determine phase velocities, liquid holdup and effective interfacial area.
The correlations below use the EOS phase densities together with viscosity, surface tension and diffusivity correlations on the bases in \cref{sec:properties}.
The hydraulic calculation requires positive properties and $0<h_L<\epsilon$, where $h_L$ and $\epsilon$ are liquid holdup and bed void fraction.
In the nine-species model a pressure-drop correlation is evaluated only to check the hydraulic operating range, and pressure is held constant in the balances.


\paragraph{Hydraulic quantities.}
The superficial velocities of each phase are calculated from molar flow rate, column area, and molar density:
\begin{align}
u^{\liq} &= \frac{\molarflow^{\liq}_{T}}{A \,\rho^{\liq}_{ \text{mol}}}
&& \left(\unit{m\,s^{-1}}\right)\\
u^{\vap} &= \frac{\molarflow^{\vap}_{T}}{A \,\rho^{\vap}_{ \text{mol}}}
&& \left(\unit{m\,s^{-1}}\right).
\end{align}


The effective interfacial area converts packed-bed geometry and liquid wetting into available transfer area per bed volume.
It depends on liquid density, surface tension, superficial velocity, and packing geometry.
The adopted expression follows the form reported by Tsai~\cite{Tsai2010}; \cref{tab:transport-applicability} compares its coefficients and source range.
\begin{align}
a_e
&= 1.42\,a_p\left[\frac{\rho^{\liq}}{\sigma^{\liq}}\,g^{1/3}\left(\frac{u^{\liq}\epsilon}{a_p}\right)^{4/3}\right]^{0.12} && \left(\unit{m^{-1}}\right)
\end{align}

Here \(\rho^{\liq}\) is liquid mass density, \(a_p\) is packing specific area and \(g\) is gravitational acceleration.


Liquid holdup, the liquid volume per packed-bed volume, enters the transfer-coefficient correlations of Chinen et al.\ and Tsai~\cite{Chinen2018,Tsai2010}:
\begin{align}
h_{L,\mathrm{raw}} &= 11.4474\left[3.185966\,u^{\liq}\left(\frac{\mu^{\liq}}{\rho_{\mathrm{mass}}^{\liq}}\right)^{1/3}\right]^{0.6471}
&& \text{dimensionless}\\
h_L &= \frac{\epsilon h_{L,\mathrm{raw}}}{\epsilon+h_{L,\mathrm{raw}}},&
h_V &= \epsilon-h_L.
\end{align}
The empirical coefficients use velocity in \unit{m\,s^{-1}}, viscosity in \unit{Pa\,s} and mass density in \unit{kg\,m^{-3}}.
The bounded mapping is applied to the raw correlation and ensures \(0<h_L<\epsilon\) for positive inputs.

\paragraph{Mass and heat transfer.}
% M9: mass and heat transfer
Gas- and liquid-side coefficients provide the physical resistances to interphase transfer.
Their correlations depend on phase density, viscosity, diffusivity, velocity and packing geometry.
The coefficients and constants follow regression-based packed-column correlations reported by Billet and Schultes and by Chinen and co-workers~\cite{Billet1999,Chinen2018}.
\begin{align}
k^{\liq}_i &= C_L
\left(\frac{g\rho^{\liq}}{\mu^{\liq}}\right)^{1/6}
\left(\frac{1}{d_h}\right)^{1/2}
\left(\frac{u^{\liq}}{a_p}\right)^{1/3}
{D_i^{\liq}}^{1/2}
&& \left(\unit{m\,s^{-1}}\right)\\
k^{\vap}_i &= \frac{C^{\vap}}{R\,T^{\vap}}
\left(\frac{a_p}{d_h\,h_V}\right)^{1/2}
\left(\frac{u^{\vap}\rho^{\vap}}{a_p\,\mu^{\vap}}\right)^{3 / 4}
{D_i^{\vap}}^{2 / 3}
\left(\frac{\mu^{\vap}}{\rho^{\vap}}\right)^{1 / 3}
&& \left(\unit{mol\,Pa^{-1}\,m^{-2}\,s^{-1}}\right)
\end{align}

Here \(d_h=4\epsilon/a_p\) is the hydraulic diameter; \(\rho^{\liq}\) and \(\rho^{\vap}\) are phase mass densities.
The adopted packing coefficients are \(C_L=0.203\) and \(C^{\vap}=0.35\).


The overall heat-transfer coefficient is estimated with the Chilton--Colburn analogy:
\begin{align}
U_T = \left(\frac{(k^{\vap}_{\COtwo}\,P)^3\,{k_{t}^{\vap}}^{2}\,C^{\vap}_{p,T}}{\left(\rho^{\vap}_{\text{mol}}D^{\vap}_{\COtwo}\right)^2}\right)^{\frac{1}{3}}
&& \left(\unit{J\,s^{-1}\,m^{-2}\,K^{-1}}\right)
\end{align}

\paragraph{Reaction enhancement.}
% M10: enhancement and film approximation
A conventional enhancement factor represents reaction-enhanced liquid-side \COtwo{} transfer as a ratio to physical absorption.
The adopted pseudo-second-order approximation follows the form of Gaspar and Fosb\o l \cite{Gaspar2015}, using bulk free-species concentrations without solving coupled reaction--diffusion equations in the liquid film.
The rate expression and Hatta number are
\begin{align}
k_{\MEA} &= 2.003\times10^4\exp(-4742.0/T^{\liq})
&& \left(\unit{m^6\,mol^{-2}\,s^{-1}}\right)\\
k_{\HtwoO} &= 4.147\exp(-3110/T^{\liq})
&& \left(\unit{m^6\,mol^{-2}\,s^{-1}}\right)\\
k_2 &= k_{\MEA}C^{\liq}_{\MEA,true}+k_{\HtwoO}C^{\liq}_{\HtwoO,true}
&& \left(\unit{m^3\,mol^{-1}\,s^{-1}}\right)\\
Ha &= \frac{\sqrt{k_2 C^{\liq}_{\MEA,true} D^{\liq}_{\COtwo}}}{k^{\liq}_{\COtwo}}
&& \text{dimensionless}
\end{align}
Here $k_2$ is the pseudo-second-order rate coefficient and $Ha$ compares the rate of reaction in the liquid film to diffusion through that film.
The enhancement calculation uses an effective dissolved-\COtwo{} concentration
$\widetilde C^{\liq}_{\COtwo}=C^{\liq}_{\COtwo,true}/s_{\COtwo}$, with the fixed dimensionless divisor
$s_{\COtwo}=1.0454$ (full precision in \cref{sec:selected-reactive-inputs}).
This empirical correction applies only within the enhancement calculation, not to species balances or EOS fugacity.
The dimensionless resistance terms are
\begin{align}
R_+ &= \frac{D^{\liq}_{\MEA} C^{\liq}_{\MEA,true}}
{2D^{\liq}_{\MEAH} C^{\liq}_{\MEAH,true}}
&& \text{dimensionless}\\
R_- &= \frac{D^{\liq}_{\MEA} C^{\liq}_{\MEA,true}}
{2D^{\liq}_{\MEACOO} C^{\liq}_{\MEACOO,true}}
&& \text{dimensionless}\\
\widehat{E} &= \frac{D^{\liq}_{\MEA} C^{\liq}_{\MEA,true}}
{2D^{\liq}_{\COtwo} \widetilde C^{\liq}_{\COtwo}}
&& \text{dimensionless}
\end{align}
which gives the explicit enhancement factor
\begin{align}
E &= 1+\frac{Ha-1}{1 + Ha(R_+ + R_- + 2)/\widehat{E}}
&& \text{dimensionless}.
\end{align}
Numerical evaluation bounds the denominators away from zero and restricts $E$ to $1\leq E\leq10^4$; the raw holdup correlation is bounded at \(10^{-12}\epsilon\).

The conventional two-film approximation relates interfacial fugacity to free molecular concentration with the bulk activity coefficient held fixed across the liquid film:
\begin{equation}
H_{\mathrm{bulk}}=\frac{f^{\ell}_{\COtwo}}{C^{\ell,\mathrm{free}}_{\COtwo}},
\qquad \Psi_H=\frac{E k^{\ell}_{\COtwo}}{E k^{\ell}_{\COtwo}+k^{v}_{\COtwo}H_{\mathrm{bulk}}}.
\end{equation}
The coefficient $H_{\mathrm{bulk}}$, the ePC-SAFT Henry coefficient below, has units of \unit{Pa\,m^3\,mol^{-1}} and uses fugacity and free \COtwo{} concentration from the same bulk equilibrium state.
Bulk equilibrium determines the liquid state, while enhancement approximates film chemistry.
\Cref{tab:transport-applicability} gives the source ranges of these correlations; \cref{sec:results} compares them with the calculated states.

% M11: flux construction
\paragraph{Interphase sources.}
The interphase molar source per unit packed length determines how $\COtwo$ and $\HtwoO$ transfer between phases along the packed height.
\begin{align}
\molarflux^{\vap}_{ \COtwo} &=  -k^{\vap}_{ \COtwo}  a_e A  \left(f^{\:\vap}_{ \COtwo} - f^{\:\liq}_{ \COtwo}\right) \Psi _{H} && \left(\unit{mol\,s^{-1}\,m^{-1}}\right) \\
\molarflux^{\liq}_{ \COtwo} &= -\molarflux^{\vap}_{ \COtwo} && \left(\unit{mol\,s^{-1}\,m^{-1}}\right) \\
\molarflux^{\vap}_{ \HtwoO} &= -k^{\vap}_{ \HtwoO} a_e A \left(f^{\:\vap}_{ \HtwoO} - f^{\:\liq}_{ \HtwoO}\right)  && \left(\unit{mol\,s^{-1}\,m^{-1}}\right) \\
\molarflux^{\liq}_{ \HtwoO} &= -\molarflux^{\vap}_{ \HtwoO} && \left(\unit{mol\,s^{-1}\,m^{-1}}\right)
\end{align}
Here $N_i^j$ is an interphase source per unit packed length, in \unit{mol.s^{-1}.m^{-1}}.
Positive $f_i^v-f_i^\ell$ gives a negative vapor source and an equal positive liquid source.
The area $a_eA$ converts interfacial transfer per area into this axial source; changes in the phase state recalculate its driving force, coefficients and enhancement.


The interphase enthalpy source contains the enthalpy carried by transferred \COtwo{} and \HtwoO{} plus sensible heat exchange between phases.
Heat transfer is modeled with the same effective interfacial area and the overall heat-transfer coefficient $U_T$:
\begin{align}
q^{\vap} &= -U_T a_e A \left(T^{\vap}-T^{\liq}\right) && \left(\unit{J\,s^{-1}\,m^{-1}}\right) \\
q^{\liq} &= -U_T a_e A \left(T^{\liq}-T^{\vap}\right) && \left(\unit{J\,s^{-1}\,m^{-1}}\right) \\
\enthalpyflow^{\vap}_{flux} &=   \molarflux^{\vap}_{ \COtwo}H^{\vap}_{ \COtwo} + \molarflux^{\vap}_{ \HtwoO}H^{\vap}_{ \HtwoO} +  q^{\vap} && \left(\unit{J\,s^{-1}\,m^{-1}}\right) \\
\enthalpyflow^{\liq}_{flux} &= -\left(\molarflux^{\vap}_{ \COtwo}H^{\vap}_{ \COtwo} + \molarflux^{\vap}_{ \HtwoO}H^{\vap}_{ \HtwoO}\right) +  q^{\liq} && \left(\unit{J\,s^{-1}\,m^{-1}}\right)
\end{align}

Here $H^{\vap}_i$ is the partial molar enthalpy of component $i$ in the wet vapor at $T^\vap$, so transferred \COtwo{} and water carry vapor-phase enthalpy across the interface.
The liquid-side enthalpy source is written as the negative vapor-side transferred enthalpy plus heat transfer so that interphase enthalpy transfer is conservative.
The heat of absorption therefore appears through the difference between these vapor partial enthalpies and the liquid partial enthalpies in the temperature equation below.

\subsection{Column balances and boundary conditions}\label{subsec:governing-equations}
\subsubsection{Material balance}\label{subsubsec:material-balance}
% M12: material conservation
Phase component flows change through the interfacial sources with signs set by the countercurrent coordinate.
Only \COtwo{} and \HtwoO{} transfer between phases; their component balances are
\begin{align}
\odv{\molarflow^{\liq}_{ \COtwo}}{z} &= -\molarflux^{\liq}_{ \COtwo} + \delta_F && \left(\unit{mol\,s^{-1}\,m^{-1}}\right)\\
\odv{\molarflow^{\liq}_{ \HtwoO}}{z} &= -\molarflux^{\liq}_{ \HtwoO} + \delta_F && \left(\unit{mol\,s^{-1}\,m^{-1}}\right)\\
\odv{\molarflow^{\vap}_{ \COtwo}}{z} &= \molarflux^{\vap}_{ \COtwo} + \delta_F && \left(\unit{mol\,s^{-1}\,m^{-1}}\right)\\
\odv{\molarflow^{\vap}_{ \HtwoO}}{z} &= \molarflux^{\vap}_{ \HtwoO} + \delta_F && \left(\unit{mol\,s^{-1}\,m^{-1}}\right)\label{eq:material-balance}
\end{align}
The evaluated flow equations include the common numerical offset \(\delta_F=\qty{1e-10}{mol\,s^{-1}\,m^{-1}}\).
It is retained in the temperature chain rule and cancels in the signed countercurrent component balance, \(\odv{(\molarflow_i^\vap-\molarflow_i^\liq)}{z}=0\).
MEA, \Ntwo, and \Otwo{} are treated as non-transferring apparent components over the one-bed calculation.
Phase compositions and loading are reconstructed from the solved flow variables:
\begin{align}
x_i &= \frac{\molarflow^{\liq}_{ i}}{\molarflow^{\liq}_{ T}}
&& \text{dimensionless}\\
y_i &= \frac{\molarflow^{\vap}_{ i}}{\molarflow^{\vap}_{ T}}
&& \text{dimensionless}\\
\alpha &= \frac{x_{\COtwo}}{x_{\MEA}}
&& \text{dimensionless}
\end{align}

\subsubsection{Energy balance}\label{subsubsec:energy-balance}
% M13: continuous energy equations
Nonisothermal balances use the ePC-SAFT phase enthalpies and the heat-transfer expression above.
Interphase heat and material transfer supply the enthalpy-flow derivatives:
\begin{align}
\odv{\enthalpyflow^{\liq}}{z} &= -\enthalpyflow^{\liq}_{flux} && \left(\unit{J\,s^{-1}\,m^{-1}}\right)\\
\odv{\enthalpyflow^{\vap}}{z} &= \enthalpyflow^{\vap}_{flux} && \left(\unit{J\,s^{-1}\,m^{-1}}\right)\label{eq:energy-balance}
\end{align}
The liquid sign follows its flow opposite increasing $z$, so the continuous equations give $\odv{(\enthalpyflow^{\vap}-\enthalpyflow^{\liq})}{z}=0$.
The liquid enthalpy flow is the enthalpy of the reacting nine-species liquid at fixed apparent \COtwo{}, MEA and water amounts, so speciation shifts with temperature and composition are included.
The vapor enthalpy flow is that of the wet four-species vapor.
With apparent partial molar enthalpies $H_i^j=(\partial\enthalpyflow^j/\partial\molarflow_i^j)_{T,P}$ of the phase-$j$ components, the temperature chain rule is
\begin{align}
\odv{\temperature^j}{z} &= \frac{\odv{\enthalpyflow^j}{z} - \sum_i H_i^j\odv{\molarflow_i^j}{z}}{\molarflow^j \left(\pdv{H^j}{\temperature^j}\right)_{P,\boldsymbol{x}^j}}
&& \left(\unit{K\,m^{-1}}\right)
\end{align}
For the transferring components \(\mathcal T=\{\COtwo,\HtwoO\}\), substitution gives
\begin{align}
\odv{\temperature^\liq}{z}
&=\frac{-\enthalpyflow^\liq_{flux}+\sum_{i\in\mathcal T}H_i^\liq(\molarflux_i^\liq-\delta_F)}
{\molarflow^\liq(\partial H_T^\liq/\partial T^\liq)_{\boldsymbol{x}}}
&&\left(\unit{K\,m^{-1}}\right),\\
\odv{\temperature^\vap}{z}
&=\frac{\enthalpyflow^\vap_{flux}-\sum_{i\in\mathcal T}H_i^\vap(\molarflux_i^\vap+\delta_F)}
{\molarflow^\vap C_{p,T}^\vap}
&&\left(\unit{K\,m^{-1}}\right).
\end{align}
At fixed apparent composition, the vapor derivative is the wet-gas heat capacity $C_{p,T}^{\vap}$ and the liquid derivative is the heat capacity of the reacting liquid, including the enthalpy of the speciation shift with temperature.
The ePC-SAFT calculation supplies both enthalpies, their partial enthalpies and their temperature derivatives from the residual Helmholtz energy and the ideal-gas records.
Numerical variation of the net enthalpy flow is measured separately in \cref{eq:energy-range}.

\subsubsection{State vector and pressure}
% M14: differential states
The differential state comprises four transferring component flows, two phase temperatures and pressure.
The nine-species model holds pressure constant over the bed:
\begin{align}
\odv{\pressure}{z} = 0
&& \left(\unit{Pa\,m^{-1}}\right).
\end{align}

The differential state vector used for the coupled calculations and its derivative mapping are given in \cref{eq:state-vector,eq:state-derivative}:
\begin{align}
\mathbf{y}(z) &=
\begin{pmatrix}
\molarflow^{\liq}_{\COtwo}(z)\\[4pt]
\molarflow^{\liq}_{\HtwoO}(z)\\[4pt]
\molarflow^{\vap}_{\COtwo}(z)\\[4pt]
\molarflow^{\vap}_{\HtwoO}(z)\\[4pt]
\temperature^{\liq}(z)\\[4pt]
\temperature^{\vap}(z)\\[4pt]
\pressure(z)
\end{pmatrix}
&& \left(
4\times\unit{mol\,s^{-1}},
2\times\unit{K},
\unit{Pa}
\right)^{\mathsf{T}}\label{eq:state-vector},\\[1em]
\odv{\mathbf{y}}{z} &=
\begin{pmatrix}
\odv{\molarflow^{\liq}_{ \COtwo}}{z} \\[4pt]
\odv{\molarflow^{\liq}_{ \HtwoO}}{z} \\[4pt]
\odv{\molarflow^{\vap}_{ \COtwo}}{z} \\[4pt]
\odv{\molarflow^{\vap}_{ \HtwoO}}{z} \\[4pt]
\odv{\temperature^{\liq}}{z} \\[4pt]
\odv{\temperature^{\vap}}{z}\\[4pt]
\odv{\pressure}{z}
\end{pmatrix}
\equiv
\mathbf{f}(z, \mathbf{y})
&& \left(
4\times\unit{mol\,s^{-1}\,m^{-1}},
2\times\unit{K\,m^{-1}},
\unit{Pa\,m^{-1}}
\right)^{\mathsf{T}}.\label{eq:state-derivative}
\end{align}
The dimensional coordinate is $z\in[0,H]$.
For normalized height $\zeta=z/H$, the solver uses $d\mathbf y/d\zeta=H\mathbf f$; figures display this normalized coordinate.

\subsubsection{Boundary conditions}
% M15: mixed boundaries
Gas-inlet and liquid-inlet specifications define mixed boundary conditions because the phases enter at opposite ends.
Let $a=0$ denote the vapor inlet and $b=H$ denote the liquid inlet.
The temperature-state calculation requires seven inlet boundary conditions:
\begin{align}
\boldsymbol{\alpha} &= \left[\molarflow^{\vap}_{ \COtwo, a}, \molarflow^{\vap}_{ \HtwoO, a}, \temperature^{\vap}_{ a}, \pressure_{a}\right]
&& \left(
2\times\unit{mol\,s^{-1}},
\unit{K},
\unit{Pa}
\right)^{\mathsf{T}}\label{eq:boundary-conditions}\\
\boldsymbol{\beta} &= \left[\molarflow^{\liq}_{ \COtwo, b}, \molarflow^{\liq}_{ \HtwoO, b}, \temperature^{\liq}_{ b}\right]
&& \left(
2\times\unit{mol\,s^{-1}},
\unit{K}
\right)^{\mathsf{T}}
\end{align}
or equivalently, the boundary residual is
\begin{align}
\mathbf{r}(\mathbf{y}) &=
\begin{pmatrix}
\mathbf{y}_{\vap}(a)-\boldsymbol{\alpha}\\
\mathbf{y}_{\liq}(b)-\boldsymbol{\beta}
\end{pmatrix}
=
\mathbf{0}
&& \left(
2\times\unit{mol\,s^{-1}},
\unit{K},
\unit{Pa},
2\times\unit{mol\,s^{-1}},
\unit{K}
\right)^{\mathsf{T}}.\label{eq:boundary-residual}
\end{align}

The specified gas flows are reconstructed from reported dry inputs and water vapor saturated at the gas-inlet temperature and pressure, while liquid flows follow reported mass flow, solvent composition and lean loading.
Saturation solves equality of the vapor water fugacity with that of pure liquid water at fixed dry-gas proportions, with both phases evaluated by the same PC-SAFT water parameters.
MEA and inert-gas flows remain constant, and the outlet transferring flows and phase temperatures are calculated rather than imposed.

\subsection{Twelve-state conserved model}\label{subsec:twelve-state}
% M16: twelve-state states and balances
The twelve-state model writes the same countercurrent bed as conservative cell balances solved simultaneously at all axial nodes.
Its states at each node are four apparent component flows, the liquid and gas temperatures, pressure, liquid holdup, the \COtwo{} and water interphase fluxes, the total energy flux and the interface log-loading.
Seven conserved balances and five algebraic relations close these states.
Seven boundary equations specify the gas inlet at the bottom and the liquid inlet at the top.
The phase enthalpies are the ePC-SAFT total enthalpies, the liquid density is the density of the solved reactive state, and pressure falls along the bed through a hydraulic pressure-drop correlation.

% M17: twelve-state film descriptions
Three descriptions of the liquid film are compared within this model.
The finite-rate film combines the film conductance along the equilibrium manifold at each axial position with a kinetic resistance from a Hatta-number film solution, using the rate expression of \cref{subsec:transport}.
Its interfacial \COtwo{} solubility is either the ePC-SAFT Henry coefficient $H_{\mathrm{bulk}}$ or the \NtwoO{}-analogy Henry coefficient of \cref{sec:properties}, which gives four film variants.
The equilibrium-manifold film assumes chemical equilibrium throughout the film, and the enhancement-factor film uses the conventional enhancement factor of \cref{subsec:transport}.
All three use nine-point film quadrature.

% M18: twelve-state input basis
The twelve-state inputs differ from the nine-species inputs.
They read the reported gas mass flow as wet gas, set the inlet water content from a pure-water vapor-pressure correlation and start from the reported gas-inlet pressure.
In Case 3C the twelve-state dry-gas flow is therefore about \qty{5}{\percent} lower than the nine-species value.
Its gas phase uses a separate four-species PC-SAFT parameter set, with a \COtwo{} dispersion energy of \qty{173.44}{K} and a constant \COtwo--water $k_{ij}$ of \num{0.01326}, instead of the published parameters of the nine-species vapor.
Capture differences between the two models therefore combine differences in model structure and inputs.
